% =========================================================
% Universal, Existential, Sentence, Open Formula, and Scope
% =========================================================

\subsection*{Basic Syntax and Semantics}

\begin{remark*}[Section toolkit: Quantifier grammar]
This section records the basic quantifier forms and the boundary between open
formulas and sentences. Quantifiers bind variables. A formula with free
variables requires an assignment for semantic evaluation; a sentence has no
free variables and has a truth value in a structure.
\end{remark*}

\begin{remark*}[Exposition]
Sentences have no free variables and have truth values in structures. Open
formulas may contain free variables and require assignments. Quantifiers turn
some free occurrences into bound occurrences by placing them inside a scope.
\end{remark*}

\begin{definitionbox}{Definition (Universal Formula)}
\begin{definition}[Universal Formula]
\label{def:universal-q}
A \emph{universal formula} is a well-formed formula of the form
\[
\forall x\,\varphi,
\]
where \(x\) is a variable and \(\varphi\) is a well-formed formula.
\end{definition}
\end{definitionbox}

\begin{remark*}[Standard quantified statement]
\[
\forall x\,\varphi\in\mathsf{Form}_{\mathcal L}
\Longleftrightarrow
x\in\mathsf{Var}_{\mathcal L}
\text{ and }
\varphi\in\mathsf{Form}_{\mathcal L}.
\]
\end{remark*}

\begin{remark*}[Definition predicate reading]
\(\operatorname{UniversalFormula}(\Phi)\) means that \(\Phi\) is a formula
whose outermost constructor is a universal quantifier.
\end{remark*}

\begin{remark*}[Negated quantified statement]
\[
\Phi \text{ is not of the form } \forall x\,\varphi.
\]
\end{remark*}

\begin{remark*}[Interpretation]
A universal formula asserts that its scope holds for every value assigned to
the bound variable.
\end{remark*}

\begin{remark*}[Examples]
\(\forall x\,P(x)\) and \(\forall x\,\exists y\,R(x,y)\) are universal
formulas.
\end{remark*}

\begin{remark*}[Non-Examples]
\(P(x)\) and \(\exists x\,P(x)\) are not universal formulas.
\end{remark*}

\begin{dependencies}
\begin{itemize}
  \item \hyperref[def:variable]{Variable}
\end{itemize}
\end{dependencies}

\begin{definitionbox}{Definition (Existential Formula)}
\begin{definition}[Existential Formula]
\label{def:existential-q}
An \emph{existential formula} is a well-formed formula of the form
\[
\exists x\,\varphi,
\]
where \(x\) is a variable and \(\varphi\) is a well-formed formula.
\end{definition}
\end{definitionbox}

\begin{remark*}[Standard quantified statement]
\[
\exists x\,\varphi\in\mathsf{Form}_{\mathcal L}
\Longleftrightarrow
x\in\mathsf{Var}_{\mathcal L}
\text{ and }
\varphi\in\mathsf{Form}_{\mathcal L}.
\]
\end{remark*}

\begin{remark*}[Definition predicate reading]
\(\operatorname{ExistentialFormula}(\Phi)\) means that \(\Phi\) is a formula
whose outermost constructor is an existential quantifier.
\end{remark*}

\begin{remark*}[Negated quantified statement]
\[
\Phi \text{ is not of the form } \exists x\,\varphi.
\]
\end{remark*}

\begin{remark*}[Interpretation]
An existential formula asserts that its scope holds for at least one value of
the bound variable.
\end{remark*}

\begin{remark*}[Examples]
\(\exists x\,P(x)\) and \(\exists x\,\forall y\,R(x,y)\) are existential
formulas.
\end{remark*}

\begin{remark*}[Non-Examples]
\(P(x)\) and \(\forall x\,P(x)\) are not existential formulas.
\end{remark*}

\begin{dependencies}
\begin{itemize}
  \item \hyperref[def:variable]{Variable}
\end{itemize}
\end{dependencies}

\begin{definitionbox}{Definition (Bounded Quantifier)}
\begin{definition}[Bounded Quantifier]
\label{def:bounded-q}
A \emph{bounded quantifier} is an abbreviation that restricts a quantifier to a
specified set or predicate:
\[
\forall x\in A\,\varphi
\quad\coloneqq\quad
\forall x\,(x\in A\to\varphi),
\]
\[
\exists x\in A\,\varphi
\quad\coloneqq\quad
\exists x\,(x\in A\land\varphi).
\]
\end{definition}
\end{definitionbox}

\begin{remark*}[Standard quantified statement]
\[
\forall x\in A\,\varphi \equiv \forall x\,(x\in A\to\varphi),
\qquad
\exists x\in A\,\varphi \equiv \exists x\,(x\in A\land\varphi).
\]
\end{remark*}

\begin{remark*}[Definition predicate reading]
\(\operatorname{BoundedQuantifier}(\Phi)\) means that \(\Phi\) uses a
restricted universal or existential quantifier.
\end{remark*}

\begin{remark*}[Negated quantified statement]
\[
\Phi \text{ is not an abbreviation of either displayed bounded form.}
\]
\end{remark*}

\begin{remark*}[Interpretation]
The restricted universal uses implication, while the restricted existential
uses conjunction. This preserves the usual behavior on empty sets.
\end{remark*}

\begin{remark*}[Examples]
\(\forall x\in A\,P(x)\) and \(\exists x\in A\,P(x)\) are bounded-quantifier
formulas.
\end{remark*}

\begin{remark*}[Non-Examples]
\(\forall x\,(x\in A\land P(x))\) is not the correct expansion of the bounded
universal form.
\end{remark*}

\begin{dependencies}
\begin{itemize}
  \item \hyperref[def:universal-q]{Universal Formula}
  \item \hyperref[def:existential-q]{Existential Formula}
\end{itemize}
\end{dependencies}

\begin{definitionbox}{Definition (Sentence)}
\begin{definition}[Sentence]
\label{def:sentence-predicate-logic}
A \emph{sentence} is a well-formed formula with no free variables.
\end{definition}
\end{definitionbox}

\begin{remark*}[Standard quantified statement]
\[
\varphi \text{ is a sentence}
\Longleftrightarrow
\varphi\in\mathsf{Form}_{\mathcal L}
\text{ and }
\mathsf{FV}(\varphi)=\varnothing.
\]
\end{remark*}

\begin{remark*}[Definition predicate reading]
\(\operatorname{Sentence}(\varphi)\) means that \(\varphi\) is a well-formed
formula with no free variables.
\end{remark*}

\begin{remark*}[Negated quantified statement]
\[
\varphi\notin\mathsf{Form}_{\mathcal L}
\quad\text{or}\quad
\mathsf{FV}(\varphi)\neq\varnothing.
\]
\end{remark*}

\begin{remark*}[Interpretation]
Sentences have truth values in structures without reference to a particular
assignment.
\end{remark*}

\begin{remark*}[Examples]
\(\forall x\,P(x)\) and \(\exists x\,\forall y\,R(x,y)\) are sentences when all
variables shown are bound.
\end{remark*}

\begin{remark*}[Non-Examples]
\(P(x)\) and \(\exists y\,R(x,y)\) are not sentences when \(x\) remains free.
\end{remark*}

\begin{dependencies}
\begin{itemize}
  \item \hyperref[def:wff-pred]{Well-Formed Formula}
\end{itemize}
\end{dependencies}

\begin{definitionbox}{Definition (Open Formula)}
\begin{definition}[Open Formula]
\label{def:open-formula-quantifiers}
An \emph{open formula} is a well-formed formula with at least one free
variable.
\end{definition}
\end{definitionbox}

\begin{remark*}[Standard quantified statement]
\[
\varphi \text{ is open}
\Longleftrightarrow
\varphi\in\mathsf{Form}_{\mathcal L}
\text{ and }
\mathsf{FV}(\varphi)\neq\varnothing.
\]
\end{remark*}

\begin{remark*}[Definition predicate reading]
\(\operatorname{OpenFormula}(\varphi)\) means that \(\varphi\) has at least one
free variable.
\end{remark*}

\begin{remark*}[Negated quantified statement]
\[
\varphi\notin\mathsf{Form}_{\mathcal L}
\quad\text{or}\quad
\mathsf{FV}(\varphi)=\varnothing.
\]
\end{remark*}

\begin{remark*}[Interpretation]
Open formulas may become true or false only after a structure and assignment
are fixed.
\end{remark*}

\begin{remark*}[Examples]
\(P(x)\) and \(R(x,y)\) are open formulas.
\end{remark*}

\begin{remark*}[Non-Examples]
\(\forall x\,P(x)\) is a sentence rather than an open formula.
\end{remark*}

\begin{dependencies}
\begin{itemize}
  \item \hyperref[def:wff-pred]{Well-Formed Formula}
  \item \hyperref[def:predicate]{Open Formula}
\end{itemize}
\end{dependencies}

\begin{definitionbox}{Definition (Quantifier Scope)}
\begin{definition}[Quantifier Scope]
\label{def:quantifier-scope}
In a formula \(Qx\,\varphi\), where \(Q\in\{\forall,\exists\}\), the
\emph{scope} of the quantifier \(Qx\) is the formula \(\varphi\). Occurrences
of \(x\) inside that scope are bound by the quantifier unless an inner
quantifier for \(x\) intervenes.
\end{definition}
\end{definitionbox}

\begin{remark*}[Standard quantified statement]
\[
Qx\,\varphi \text{ has scope } \varphi,
\qquad
Q\in\{\forall,\exists\}.
\]
\end{remark*}

\begin{remark*}[Definition predicate reading]
\(\operatorname{Scope}(x,\varphi)\) means that \(\varphi\) is the scope governed
by the quantifier binding \(x\).
\end{remark*}

\begin{remark*}[Negated quantified statement]
\[
\varphi \text{ is not the formula governed by the displayed quantifier on } x.
\]
\end{remark*}

\begin{remark*}[Interpretation]
Scope determines which variable occurrences a quantifier controls. Scope errors
are the main source of incorrect translations from English into predicate
logic.
\end{remark*}

\begin{remark*}[Examples]
In \(\forall x\,(P(x)\to Q(x))\), the scope is \((P(x)\to Q(x))\). In
\((\forall x\,P(x))\to Q(x)\), the final occurrence of \(x\) is outside the
scope of the quantifier.
\end{remark*}

\begin{dependencies}
\begin{itemize}
  \item \hyperref[def:universal-q]{Universal Formula}
  \item \hyperref[def:existential-q]{Existential Formula}
\end{itemize}
\end{dependencies}

\begin{theorembox}{Theorem (Negation of Bounded Quantifiers)}
\begin{theorem}[Negation of Bounded Quantifiers]
\label{thm:neg-bounded}
Let \(A\) be a set and let \(\varphi\) be a formula. Then
\[
\neg(\forall x\in A\,\varphi)\equiv \exists x\in A\,\neg\varphi,
\qquad
\neg(\exists x\in A\,\varphi)\equiv \forall x\in A\,\neg\varphi.
\]
\hyperref[prf:neg-bounded]{Go to proof.}
\end{theorem}
\end{theorembox}

\begin{remark*}[Standard quantified statement]
\[
\neg\forall x\,(x\in A\to\varphi)
\equiv
\exists x\,(x\in A\land\neg\varphi),
\qquad
\neg\exists x\,(x\in A\land\varphi)
\equiv
\forall x\,(x\in A\to\neg\varphi).
\]
\end{remark*}

\begin{remark*}[Predicate reading]
The theorem records the negation law for bounded universal and bounded
existential formulas.
\end{remark*}

\begin{remark*}[Interpretation]
The domain restriction is preserved under negation. What changes is the
quantifier and the inner assertion.
\end{remark*}

\begin{dependencies}
\begin{itemize}
  \item \hyperref[def:bounded-q]{Bounded Quantifier}
  \item \hyperref[thm:qneg]{Quantifier Negation Laws}
\end{itemize}
\end{dependencies}
